% Options for packages loaded elsewhere
\PassOptionsToPackage{unicode}{hyperref}
\PassOptionsToPackage{hyphens}{url}
\PassOptionsToPackage{dvipsnames,svgnames,x11names}{xcolor}
\PassOptionsToPackage{space}{xeCJK}
\documentclass[
  11pt,
  a4paper,
]{article}
\usepackage{xcolor}
\usepackage[margin=1in]{geometry}
\usepackage{amsmath,amssymb}
\setcounter{secnumdepth}{-\maxdimen} % remove section numbering
\usepackage{iftex}
\ifPDFTeX
  \usepackage[T1]{fontenc}
  \usepackage[utf8]{inputenc}
  \usepackage{textcomp} % provide euro and other symbols
\else % if luatex or xetex
  \usepackage{unicode-math} % this also loads fontspec
  \defaultfontfeatures{Scale=MatchLowercase}
  \defaultfontfeatures[\rmfamily]{Ligatures=TeX,Scale=1}
\fi
\usepackage{lmodern}
\ifPDFTeX\else
  % xetex/luatex font selection
  \setmainfont[]{Times New Roman}
  \setmonofont[]{Consolas}
  \ifXeTeX
    \usepackage{xeCJK}
    \setCJKmainfont[]{Microsoft YaHei}
    \setCJKmonofont[]{Microsoft YaHei}
  \fi
  \ifLuaTeX
    \usepackage[]{luatexja-fontspec}
    \setmainjfont[]{Microsoft YaHei}
  \fi
\fi
% Use upquote if available, for straight quotes in verbatim environments
\IfFileExists{upquote.sty}{\usepackage{upquote}}{}
\IfFileExists{microtype.sty}{% use microtype if available
  \usepackage[]{microtype}
  \UseMicrotypeSet[protrusion]{basicmath} % disable protrusion for tt fonts
}{}
\makeatletter
\@ifundefined{KOMAClassName}{% if non-KOMA class
  \IfFileExists{parskip.sty}{%
    \usepackage{parskip}
  }{% else
    \setlength{\parindent}{0pt}
    \setlength{\parskip}{6pt plus 2pt minus 1pt}}
}{% if KOMA class
  \KOMAoptions{parskip=half}}
\makeatother
\setlength{\emergencystretch}{3em} % prevent overfull lines
\providecommand{\tightlist}{%
  \setlength{\itemsep}{0pt}\setlength{\parskip}{0pt}}
\providecommand{\bm}[1]{\symbfit{#1}}
\usepackage{bookmark}
\IfFileExists{xurl.sty}{\usepackage{xurl}}{} % add URL line breaks if available
\urlstyle{same}
\hypersetup{
  colorlinks=true,
  linkcolor={Maroon},
  filecolor={Maroon},
  citecolor={Blue},
  urlcolor={Blue},
  pdfcreator={LaTeX via pandoc}}

\author{}
\date{}

\begin{document}

\section{第02讲：统计语言模型}\label{ux7b2c02ux8bb2ux7edfux8ba1ux8bedux8a00ux6a21ux578b}

\textbf{DATA130030.01，自然语言处理}（阅读时长约60分钟）\\
\textbf{日期：} 2025年9月17日\\
\textbf{主讲人：} 周宝健\\
\textbf{单位：} 复旦大学数据科学学院

\begin{center}\rule{0.5\linewidth}{0.5pt}\end{center}

\begin{quote}
\textbf{译注（勘误）：} 原文开头有一处孤立的
\texttt{\$\$x\$\$}，系排版残留，已删除。原文 bigram 推导中将 \(v_0\)
误称为 ``EOS 起始标记''，应为 BOS，译文已更正。unigram 推导中对
\(\lambda\) 求偏导的求和上限原文写为 \(|\mathcal{V}|\)，应为
\(|\mathcal{V}|+1\)（含 EOS），译文已更正。
\end{quote}

本讲介绍 \(N\)-gram 语言模型。我们将定义 \(N\)-gram
语言模型，并说明如何使用最大似然估计（MLE）来学习模型参数；接着讨论处理零概率问题的技术（平滑），并定义困惑度（perplexity）。

\textbf{符号说明}\\
经过分词后，一个句子表示为 \([w_1, w_2, \ldots, w_t]\)，记作
\(s = w_{1:t}\)。例如，句子 \emph{``the cat saw the dog''} 对应
\(s = w_{1:5} = [\textit{the, cat, saw, the, dog}]\)。设 \(P(X)\)
表示离散随机变量 \(X\) 的概率分布，用 \(w_i\) 表示句子第 \(i\)
个位置的随机变量 \(X_i\)，即 \(P(X_i = w_i) \equiv P(w_i)\)。设
\(X_{i+1}\) 表示选择 \(w_{i+1}\)，\(X_i\) 表示在 \(w_{i+1}\) 之前出现的
\(w_i\)，则 \(P(X_i, X_{i+1})\) 为 \(w_i w_{i+1}\) 的联合概率，记作
\(P(X_i = w_i, X_{i+1} = w_{i+1}) \equiv P(w_{i:i+1})\)。我们定义词表
\(\mathcal{V} = \{v_1, v_2, \ldots, v_{|\mathcal{V}|}\}\)，以及两个特殊符号
\(v_0 = \text{BOS}\)（句首）和
\(v_{|\mathcal{V}|+1} = \text{EOS}\)（句尾）。

\begin{center}\rule{0.5\linewidth}{0.5pt}\end{center}

\subsection{动机}\label{ux52a8ux673a}

假设你正在做一个语音识别任务，需要设计一个语音识别器。给定一段声学信号，识别器产生两个候选句子：\(s_1 =\)
\emph{``It's hard to recognize speech''}（语音识别很难）和 \(s_2 =\)
\emph{``It's hard to wreck a nice
beach''}（毁掉一片好海滩很难）。如果模型足够好，它应使
\(P(s_1) > P(s_2)\)，因为两者的声学信号非常相似，但 \(s_1\)
在日常语言中更为合理。因此，语音识别器需要能够为不同的句子序列赋予概率，以比较
\(s_1\) 和 \(s_2\)
在现实中出现的可能性。\textbf{语言建模的任务就是为句子赋予概率}，即我们希望学到一个分布
\(P\)，使得 \(P(s_1) > P(s_2)\)。

\textbf{我们的任务：} 目标是在给定某语言词表 \(\mathcal{V}\)
下的\textbf{训练语料库}（一组句子）上学习一个语言模型。利用训练语料库学习模型后，使用\textbf{验证语料库}调整模型超参数，最后用\textbf{测试语料库}评估训练好的语言模型的性能。

\begin{center}\rule{0.5\linewidth}{0.5pt}\end{center}

\subsection{\texorpdfstring{统计语言模型与
\(N\)-gram}{统计语言模型与 N-gram}}\label{ux7edfux8ba1ux8bedux8a00ux6a21ux578bux4e0e-n-gram}

统计语言模型定义了在给定词表 \(\mathcal{V}\)
上句子的概率分布。英语语言模型的词表可能是
\(\mathcal{V} = \{\textit{the}, \textit{dog}, \textit{laughs}, \textit{saw}, \ldots\}\)，其中
\(\mathcal{V}\) 是有限集。语言中的\textbf{句子}是词的序列

\[s = w_1 w_2 \ldots w_t \equiv w_{1:t},\]

其中整数 \(t \geq 1\)
不固定，\(w_i \in \mathcal{V}\)（\(i \in \{1, \ldots, t-1\}\)），并假设
\(w_t = \text{EOS}\) 是不在 \(\mathcal{V}\) 中的特殊符号。定义
\(\mathcal{V}^{\dagger}\) 为词表 \(\mathcal{V}\)
上所有句子的集合，它是无限集（句子长度任意）。例如：

\[\mathcal{V}^{\dagger} = \{\text{the dog barks EOS},\ \text{the EOS},\ \text{cat cat cat EOS},\ \text{EOS},\ \ldots\}.\]

语言模型计算 \(w_{1:t}\) 的概率 \(p(w_{1:t})\)，其正式定义如下：

\begin{quote}
\textbf{定义（语言模型，LM）}\\
给定有限词表 \(\mathcal{V}\)，对任意 \(t \geq 1\)，长度为 \(t\)
的句子是序列 \(w_{1:t}\)，其中
\(w_i \in \mathcal{V}\)（\(i = 1, \ldots, t-1\)），\(w_t = \text{EOS}\)。设
\(\mathcal{V}^{\dagger}\)
为所有此类句子的集合。\textbf{语言模型}是定义在
\(\mathcal{V}^{\dagger}\) 上的概率分布函数 \(p(w_{1:t})\)，满足：

\begin{itemize}
\tightlist
\item
  对任意 \(w_{1:t} \in \mathcal{V}^{\dagger}\)，\(p(w_{1:t}) \geq 0\)；
\item
  \(\sum_{w_{1:t} \in \mathcal{V}^{\dagger}} p(w_{1:t}) = 1\)。
\end{itemize}
\end{quote}

由此，\(p(w_{1:t})\) 是定义在 \(\mathcal{V}^{\dagger}\)
上的概率分布。语言模型使我们能够\textbf{理解}（通过 \(p\)
比较任意两个句子的概率）和\textbf{生成}（用分布 \(p\)
生成句子）。利用链式法则，可将其分解为条件概率之积：

\[p(w_{1:t}) = \prod_{i=1}^t p(w_i \mid w_{1:i-1}),\]

其中 \(p(w_1 \mid w_{1:0} = \text{BOS}) = p(w_1)\)，\(w_{1:i-1}\) 称为
\(w_i\) 的\textbf{历史}。精确计算 \(p(w_{1:t})\)
通常极为困难。例如，仅对5个词的句子 \(p(w_{1:5})\)，在
\(|\mathcal{V}| = 10^4\) 时就需要约 \(10^{20}\)
个参数。目前所有语言模型本质上都是对 \(p(w_{1:t})\)
的\textbf{估计}，有三大类：

\begin{enumerate}
\def\labelenumi{\arabic{enumi}.}
\tightlist
\item
  \textbf{统计语言模型}：使用统计方法，基于词序列的概率预测下一个词，即
  \(N\)-gram 模型。
\item
  \textbf{神经语言模型}：基于神经网络，包括两层神经网络（word2vec）和
  RNN（如 LSTM、GRU）等架构。
\item
  \textbf{大语言模型（LLM）}：使用自注意力机制（Transformer）的神经网络模型，如
  BERT 和 GPT 系列。
\end{enumerate}

\textbf{为何需要 \(w_t = \text{EOS}\)？} 句子长度 \(t\)
本身是随机变量。最常见的处理方式是令最后一个词 \(w_t\) 始终等于特殊符号
EOS，且该符号只能出现在序列末尾。

近似 \(p(w_{1:t})\)
是构建实用语言模型的关键。在\textbf{一阶马尔可夫假设}下，\(p(w_{1:t}) = \prod_{i=1}^t P(X_i = w_i \mid X_{i-1} = w_{i-1})\)。类似地，在\textbf{二阶马尔可夫假设}下：

\[p(w_{1:t}) = \prod_{i=1}^t P(X_i = w_i \mid X_{i-2} = w_{i-2},\ X_{i-1} = w_{i-1}),\]

生成句子的步骤如下：

\begin{itemize}
\tightlist
\item
  \textbf{步骤1：} 初始化 \(i = 1\)，令 \(x_0 = x_{-1} = \text{BOS}\)；
\item
  \textbf{步骤2：} 从分布
  \(p(X_i = x_i \mid X_{i-2:i-1} = x_{i-2:i-1})\) 中采样 \(x_i\)；
\item
  \textbf{步骤3：} 若 \(x_i = \text{EOS}\)，则返回 \(x_{1:i}\)；否则令
  \(i = i + 1\)，返回步骤2。
\end{itemize}

\(N\)-gram 模型的正式定义如下：

\begin{quote}
\textbf{定义（\(N\)-gram 语言模型）}\\
给定有限词表 \(\mathcal{V}\)，\(N\)-gram 语言模型在 \((N-1)\)
阶马尔可夫假设下计算句子 \(w_{1:t}\) 的概率：

\[p(w_{1:t}) = \prod_{i=1}^t q(w_i \mid w_{i-N+1:i-1}),\]

其中 \(q(x_N \mid x_{1:N-1})\) 是 \(N\)-gram \(x_{1:N}\)
的条件概率分布，\(x_i \in \mathcal{V} \cup \{\text{BOS}\}\)（\(i = 1, \ldots, N-1\)），\(x_N \in \mathcal{V} \cup \{\text{EOS}\}\)。默认规定
\(i \leq 0\) 时 \(w_i = \text{BOS}\)。
\end{quote}

具体而言：

\begin{itemize}
\tightlist
\item
  \(N = 1\)：\textbf{一元模型（Unigram）}，\(p(w_{1:t}) = \prod_{i=1}^t q(w_i)\)；
\item
  \(N = 2\)：\textbf{二元模型（Bigram）}，\(p(w_{1:t}) = \prod_{i=1}^t q(w_i \mid w_{i-1})\)；
\item
  \(N = 3\)：\textbf{三元模型（Trigram）}，\(p(w_{1:t}) = \prod_{i=1}^t q(w_i \mid w_{i-2:i-1})\)。
\end{itemize}

模型参数约有 \(|\mathcal{V}|^N\) 个。以
\(|\mathcal{V}| = 10{,}000\)、\(N = 3\) 为例，参数数量约
\(10^{12}\)，数量极为庞大。

\begin{center}\rule{0.5\linewidth}{0.5pt}\end{center}

\subsection{\texorpdfstring{从训练语料库通过最大似然估计学习
\(q(x_N \mid x_{1:N-1})\)}{从训练语料库通过最大似然估计学习 q(x\_N \textbackslash mid x\_\{1:N-1\})}}\label{ux4eceux8badux7ec3ux8bedux6599ux5e93ux901aux8fc7ux6700ux5927ux4f3cux7136ux4f30ux8ba1ux5b66ux4e60-qx_n-mid-x_1n-1}

训练数据集为：

\[\mathcal{D}_{\text{tr}} = \{s_1, s_2, \ldots, s_T\} \equiv s_{1:T}, \quad s_i = w_{1:t_i}^{(i)} \in \mathcal{V}^{\dagger}.\]

设模型由参数向量
\(\bm{\theta} = [\theta_1, \theta_2, \ldots, \theta_d]^\top\)
刻画。\textbf{最大似然估计（MLE）}
的目标是找到使观测数据联合概率最大的参数：

\[\hat{\bm{\theta}} \in \argmax_{\bm{\theta} \in \bm{\Theta}} \left\{ \mathcal{L}_T(\bm{\theta}; \mathcal{D}_{\text{tr}}) := P(s_{1:T}; \bm{\theta}) \right\}.\]

\subsubsection{一元模型的 MLE}\label{ux4e00ux5143ux6a21ux578bux7684-mle}

在一元模型假设下：

\[P(s_{1:T}) = \prod_{i=1}^{|\mathcal{V}|+1} \theta_i^{C(v_i)},\]

其中 \(C(v_i)\) 为 \(v_i\) 在训练语料中出现的次数，参数向量
\(\bm{\theta} = [\theta_1, \ldots, \theta_{|\mathcal{V}|}, \theta_{\text{EOS}}]^\top\)，约束为
\(\sum_{i=1}^{|\mathcal{V}|+1} \theta_i = 1\)。

对数似然为：

\[\argmax_{\sum_j \theta_j = 1} \sum_{i=1}^{|\mathcal{V}|+1} C(v_i) \log \theta_i.\]

引入拉格朗日乘子 \(\lambda\)，构造拉格朗日函数：

\[G(\bm{\theta}, \lambda) = \sum_{i=1}^{|\mathcal{V}|+1} C(v_i) \log \theta_i - \lambda\left(\sum_{i=1}^{|\mathcal{V}|+1} \theta_i - 1\right).\]

令偏导为零（\textbf{勘误：} 原文求和上限为 \(|\mathcal{V}|\)，应为
\(|\mathcal{V}|+1\)）：

\[\frac{\partial G}{\partial \theta_i} = \frac{C(v_i)}{\theta_i} - \lambda = 0, \quad \frac{\partial G}{\partial \lambda} = \sum_{i=1}^{|\mathcal{V}|+1} \theta_i - 1 = 0.\]

解得：

\[\theta_i^* = \frac{C(v_i)}{\sum_{i=1}^T |s_i|}.\]

即一元模型的 MLE 参数就是\textbf{频率估计}。

\subsubsection{二元模型的 MLE}\label{ux4e8cux5143ux6a21ux578bux7684-mle}

对于二元模型，令 \(\theta_{i,j+1} := q(v_{j+1} \mid v_i)\)，约束为对所有
\(i\)：\(\sum_{j=0}^{|\mathcal{V}|} \theta_{i,j+1} = 1\)。

对数似然为：

\[\log P(s_{1:T}; \bm{\Theta}) = \sum_{i=0}^{|\mathcal{V}|} \sum_{j=0}^{|\mathcal{V}|} C(v_i, v_{j+1}) \log q(v_{j+1} \mid v_i).\]

类似地引入拉格朗日乘子 \(\lambda_i\)，令梯度为零，解得：

\[\theta_{i,j+1}^* = q(v_{j+1} \mid v_i) = \frac{C(v_i, v_{j+1})}{\sum_{j=0}^{|\mathcal{V}|} C(v_i, v_{j+1})}.\]

对三元模型及更高阶，MLE 估计为：

\[q(x_N \mid x_{1:N-1}) = \frac{C(x_{1:N})}{C(x_{1:N-1})}.\]

\begin{center}\rule{0.5\linewidth}{0.5pt}\end{center}

\subsection{\texorpdfstring{\(N\)-gram
语言模型的平滑}{N-gram 语言模型的平滑}}\label{n-gram-ux8bedux8a00ux6a21ux578bux7684ux5e73ux6ed1}

\(N\)-gram
模型通常存在\textbf{数据稀疏}问题。例如，在3800万词的新闻语料上观察所有三元组后，同来源新文章中仍有约三分之一的三元组是全新未见的。对未见
\(N\)-gram 直接用 MLE
会导致零概率，因此需要平滑技术------从高频事件中''削减''一些概率质量，分配给未见事件。

\subsubsection{加法平滑（Additive
Smoothing）}\label{ux52a0ux6cd5ux5e73ux6ed1additive-smoothing}

对每个 \(N\)-gram 计数加上一个小量 \(\delta\)（通常
\(0 < \delta \leq 1\)）：

\[q_{\text{Add}}(x_N \mid x_{1:N-1}) = \frac{\delta + C(x_{1:N})}{\delta(|\mathcal{V}|+1) + \sum_{x_N \in \mathcal{V} \cup \{\text{EOS}\}} C(x_{1:N})}.\]

最简单的情形是 \(\delta = 1\)（加一平滑）。

\subsubsection{线性插值（Linear
Interpolation）}\label{ux7ebfux6027ux63d2ux503clinear-interpolation}

以三元模型为例，将三元、二元、一元估计加权混合：

\[q_{\text{Int}}(x_3 \mid x_{1:2}) = \lambda_1 \cdot q_{\text{ML}}(x_3) + \lambda_2 \cdot q_{\text{ML}}(x_3 \mid x_2) + \lambda_3 \cdot q_{\text{ML}}(x_3 \mid x_{1:2}),\]

其中 \(\sum_{i=1}^3 \lambda_i = 1\)，各 \(\lambda_i\)
从验证集学习。一种简化的''分桶''方法（Bucketing）定义为：

\[\lambda_1 = \frac{C(x_{1:2})}{C(x_{1:2}) + \gamma}, \quad \lambda_2 = (1 - \lambda_1) \cdot \frac{C(x_2)}{C(x_2) + \gamma}, \quad \lambda_3 = 1 - \lambda_1 - \lambda_2,\]

其中 \(\gamma > 0\) 为唯一超参数。可验证各 \(\lambda_i\) 为正且和为1。

\subsubsection{Katz 回退平滑（Katz
Back-off）}\label{katz-ux56deux9000ux5e73ux6ed1katz-back-off}

对二元模型，对任意 \(x_1\) 定义：

\[\mathcal{A}(x_1) = \{x_2: C(x_1 x_2) > 0\}, \quad \mathcal{B}(x_1) = \{x_2: C(x_1 x_2) = 0\}.\]

估计为：

\[q_{\text{Katz}}(x_2 \mid x_1) = \begin{cases} \dfrac{C^*(x_1, x_2)}{C(x_1)} & \text{若 } x_2 \in \mathcal{A}(x_1) \\ \alpha(x_1) \cdot \dfrac{q_{\text{ML}}(x_2)}{\sum_{x_2 \in \mathcal{B}(x_1)} q_{\text{ML}}(x_2)} & \text{若 } x_2 \in \mathcal{B}(x_1) \end{cases}\]

其中折扣计数 \(C^*(x_1, x_2) = C(x_1, x_2) - \beta\)，典型取值
\(\beta = 0.5\)。对三元模型，类似地定义 \(\mathcal{A}(u,v)\) 和
\(\mathcal{B}(u,v)\)：

\[q_{\text{Katz}}(w \mid u, v) = \begin{cases} \dfrac{C^*(u,v,w)}{C(u,v)} & \text{若 } w \in \mathcal{A}(u,v) \\ \alpha(u,v) \cdot \dfrac{q_{\text{Katz}}(w \mid v)}{\sum_{w \in \mathcal{B}(u,v)} q_{\text{Katz}}(w \mid v)} & \text{若 } w \in \mathcal{B}(u,v) \end{cases}\]

其中
\(\alpha(u,v) = 1 - \sum_{w \in \mathcal{A}(u,v)} \frac{C^*(u,v,w)}{C(u,v)}\)。

\subsubsection{Kneser-Ney 平滑（Kneser-Ney
Smoothing）}\label{kneser-ney-ux5e73ux6ed1kneser-ney-smoothing}

给定折扣参数 \(\delta \in [0,1]\)，二元 KN 平滑为：

\[q_{\text{KN}}(x_2 \mid x_1) = \frac{\max\{C(x_1 x_2) - \delta,\ 0\}}{\sum_{x \in \mathcal{V} \cup \{\text{EOS}\}} C(x_1 x)} + \lambda_{x_1} \cdot q_{\text{KN}}(x_2),\]

其中一元概率衡量词 \(x_2\) 出现在新上下文中的可能性：

\[q_{\text{KN}}(x_2) = \frac{|\{x: C(x x_2) > 0\}|}{|\{(x_1, x_2): C(x_1 x_2) > 0\}|},\]

参数 \(\lambda_{x_1}\) 为：

\[\lambda_{x_1} = \frac{\delta \cdot |\{w: C(x_1 w) > 0\}|}{\sum_v C(x_1 v)}.\]

推广到 \(N\)-gram：

\[q_{\text{KN}}(x_N \mid x_{1:N-1}) = \frac{\max\{C(x_{1:N}) - \delta,\ 0\}}{\sum_{x \in \mathcal{V} \cup \{\text{EOS}\}} C(x_{1:N-1} x)} + \lambda_{x_{1:N-1}} \cdot q_{\text{KN}}(x_N \mid x_{2:N-1}),\]

其中：

\[\lambda_{x_{1:N-1}} = \frac{\delta \cdot |\{x: C(x_{1:N-1} x) > 0\}|}{\sum_{x \in \mathcal{V} \cup \{\text{EOS}\}} C(x_{1:N-1} x)}.\]

改进版 Kneser-Ney 平滑已在 \href{https://github.com/kpu/kenlm}{kenlm}
中实现。

\subsubsection{愚蠢回退（Stupid
Back-off）}\label{ux611aux8822ux56deux9000stupid-back-off}

不带折扣的回退（非真正的概率）：

\[S(x_N \mid x_{1:N-1}) = \begin{cases} \dfrac{C(x_{1:N})}{C(x_{1:N-1})} & \text{若 } C(x_{1:N}) > 0 \\ 0.4 \cdot S(x_N \mid x_{2:N-1}) & \text{否则} \end{cases}\]

\[S(x_N) = \frac{C(x_N)}{T},\]

其中 \(T\) 为训练语料的总大小。

\begin{center}\rule{0.5\linewidth}{0.5pt}\end{center}

\subsection{\texorpdfstring{用困惑度评估
\(N\)-gram}{用困惑度评估 N-gram}}\label{ux7528ux56f0ux60d1ux5ea6ux8bc4ux4f30-n-gram}

设测试数据为句子
\(s_1, s_2, \ldots, s_m\)，一个好的模型应对这些真实句子赋予较高概率。总体质量可用：

\[\prod_{i=1}^m p(s_i)\]

来衡量。设 \(M = \sum_{i=1}^m t_i\)
为测试语料的总词数，\textbf{平均对数概率}定义为：

\[\frac{1}{\sum_{i=1}^m |s_i|} \log_2 \prod_{i=1}^m p(s_i) = \frac{1}{\sum_{i=1}^m |s_i|} \sum_{i=1}^m \log_2 p(s_i).\]

\textbf{困惑度（Perplexity）} 定义为：

\[\operatorname{PPL}(s_{1:m}) := 2^{-l}, \quad \text{其中 } l = \frac{1}{\sum_{i=1}^m |s_i|} \sum_{i=1}^m \log_2 p(s_i).\]

困惑度是正数，\textbf{值越小，语言模型对未见数据的建模能力越强}。其直觉含义为：\emph{``如果在每一步按照语言模型的概率分布随机采样词，平均需要采样多少次才能得到正确的词？''}（Neubig,
2017, 第9页）。

\textbf{示例：} 若一元模型对所有词赋予相同概率，则
\(\operatorname{PPL}(s_{1:m}) = |\mathcal{V}|\)------平均需要从全部词表中选择。若
\(|\mathcal{V}| = 1{,}000{,}000\)
且一个好的语言模型的困惑度为30，则平均只需选约30个词。

\subsubsection{熵与交叉熵}\label{ux71b5ux4e0eux4ea4ux53c9ux71b5}

困惑度本质上衡量语言模型的\textbf{熵}。给定离散随机变量
\(X\)，\textbf{熵}定义为：

\[H(p) \triangleq -\sum_{i=1}^n p(x_i) \log p(x_i).\]

分布 \(q\) 相对于分布 \(p\) 的\textbf{交叉熵}定义为：

\[H(p, q) = -\sum_{x_i \in \mathcal{X}} p(x_i) \log q(x_i).\]

在语言模型中，\(p\) 为真实分布，\(q\)
为从训练语料估计的分布。由于真实分布未知，我们无法直接计算 \(H(p)\) 或
\(H(p,q)\)。但可以利用：

\[H(p, q) = H(p) + \underbrace{\sum_{x_i \in \mathcal{X}} p(x_i) \left(\log p(x_i) - \log q(x_i)\right)}_{D_{\text{KL}}(p \| q) \geq 0} \geq H(p).\]

因此 \(H(p,q)\) 是 \(H(p)\) 的上界，\(q\) 越接近 \(p\)，\(H(p,q)\)
越接近 \(H(p)\)（等号当且仅当 \(p = q\) 时成立）。

由 \textbf{Shannon-McMillan-Breiman 定理}（Jurafsky, 2022, 公式3.49）：

\[H(p) \leq H(p, q) = \lim_{n \to \infty} -\frac{1}{n} \log q(w_{1:n}) \approx H(s_{1:m}) \triangleq -\frac{1}{|s_{1:m}|} \log p(s_{1:m}).\]

\textbf{困惑度}正是该近似量 \(H(s_{1:m})\) 的指数形式：

\[\operatorname{PPL}(s_{1:m}) = 2^{-\frac{1}{|s_{1:m}|} \log_2 p(s_{1:m})} = 2^{H(s_{1:m})}.\]

采用指数形式没有特别原因，但它使数值差异更为显著，比 \(H(W)\)
对模型性能更为敏感。

\textbf{困惑度的优势：} 1.
比熵值更直观易记（如100--200的困惑度范围，比6.64--7.64比特的熵值更具可解释性）；
2. 10\%的困惑度改善比2\%的熵减小更具感知显著性； 3.
可直接用留出集或测试数据计算，困惑度越低说明模型越接近生成观测数据的''真实''模型。

\(N\)-gram 模型的核心组件曾被用于第一代谷歌翻译（Brants et al.,
2007）。然而，作为纯统计语言模型，它已先后被基于 RNN
的模型和大语言模型（LLM）所取代。

\begin{center}\rule{0.5\linewidth}{0.5pt}\end{center}

\subsection{参考文献}\label{ux53c2ux8003ux6587ux732e}

\begin{itemize}
\tightlist
\item
  https://cs229.stanford.edu/lectures-spring2023/cs229-probability\_review.pdf
\item
  https://cs229.stanford.edu/notes2022fall/main\_notes.pdf
\item
  https://anoopsarkar.github.io/nlp-class/assets/slides/prob.pdf
\item
  https://www.cs.columbia.edu/\textasciitilde mcollins/courses/nlp2011/notes/lm.pdf
\item
  https://pages.cs.wisc.edu/\textasciitilde jerryzhu/cs838/LM.pdf
\item
  https://www3.nd.edu/\textasciitilde dchiang/teaching/nlp/2021/notes/chapter2v2.pdf
\item
  https://www.fit.vut.cz/study/phd-thesis-file/283/283.pdf
\end{itemize}

\begin{center}\rule{0.5\linewidth}{0.5pt}\end{center}

\subsection{附录}\label{ux9644ux5f55}

\textbf{Jensen 不等式：} 对凹函数 \(\varphi: \mathbb{R} \to \mathbb{R}\)
及正权 \(a_i\)：

\[\varphi\!\left(\sum_{i=1}^n \frac{x_i}{n}\right) \geq \frac{1}{n} \sum_{i=1}^n \varphi(x_i) \quad (\text{当 } a_i = 1).\]

\begin{quote}
\textbf{定理（最大熵）}\\
对随机变量 \(X\) 及其概率分布 \(p\)，当 \(p\) 为均匀分布时，\(X\)
的熵最大。
\end{quote}

\textbf{证明：}\\
令 \(\varphi(y) = -y \log y\)（规定 \(\varphi(0) = 0\)），设 \(X\) 有
\(n\) 个可能取值 \(x_1, \ldots, x_n\)，概率为
\(p(x_1), \ldots, p(x_n)\)。由于 \(\varphi(y)\) 在 \([0,1]\)
上为凹函数，由 Jensen 不等式：

\[\begin{aligned}
n \cdot \varphi\!\left(\sum_{i=1}^n \frac{y_i}{n}\right) &\geq \sum_{i=1}^n \varphi(y_i) \\
n \cdot \left(-\frac{1}{n}\right) \log\!\left(\frac{1}{n}\right) &\geq -\sum_{i=1}^n p(x_i) \log p(x_i) := H(p),
\end{aligned}\]

即 \(\log(n) \geq H(p)\)。上界在
\(p(x_i) = 1/n\)（均匀分布）时取到，证毕。

更多内容见 Conrad (2004)。

\end{document}
